% ============================================================
% 2.6 Implementacion - App Turismo Comarapa
% Requiere en el preambulo del documento principal:
%   \usepackage{float}    % para el especificador de posicion [H]
%   \usepackage{listings} % para los fragmentos de codigo
%   \usepackage{xcolor}   % colores de los fragmentos
% Nota: los fragmentos se dejaron sin tildes ni enies dentro del
% codigo porque listings no maneja UTF-8 multibyte sin configuracion
% extra. Cada fragmento tiene 20 lineas o menos.
% ============================================================

\lstdefinestyle{codigo}{
  basicstyle=\ttfamily\footnotesize,
  keywordstyle=\color{blue!70!black}\bfseries,
  commentstyle=\color{green!40!black}\itshape,
  stringstyle=\color{red!60!black},
  numbers=left,
  numberstyle=\tiny\color{gray},
  frame=single,
  breaklines=true,
  showstringspaces=false,
  tabsize=2,
  captionpos=b
}
\lstdefinelanguage{Dart}{
  morekeywords={class,final,const,return,if,else,await,async,try,catch,
    finally,null,true,false,new,extends,import,Future,void,String,bool,
    required,throw,switch,var},
  sensitive=true,
  morecomment=[l]{//},
  morecomment=[s]{/*}{*/},
  morestring=[b]',
  morestring=[b]"
}

\subsection*{2.6 Implementación}

En esta sección se describen las funcionalidades críticas que
efectivamente se construyeron en la aplicación de turismo del municipio
de Comarapa, las decisiones técnicas que las sostienen y las
dificultades encontradas durante el desarrollo junto con la forma en que
se resolvieron. Los fragmentos de código incluidos son únicamente
ilustrativos; el código completo se encuentra en el repositorio del
proyecto.

\subsubsection*{Funcionalidades críticas construidas}

\begin{enumerate}
  \item \textbf{Autenticación y control de acceso por rol.} Inicio de
  sesión con correo y contraseña mediante Supabase Auth, con dos roles
  (\texttt{administrador} y \texttt{editor}) cuya autorización se aplica
  en la propia base de datos mediante políticas RLS.
  \item \textbf{Gestión de contenido turístico (CRUD).} Alta, edición y
  desactivación de lugares, hoteles, restaurantes, eventos, actividades
  y gastronomía desde el panel de administración, incluida la selección
  de la ubicación en el mapa y la subida de imágenes a Supabase Storage.
  \item \textbf{Gestión de usuarios (solo administrador).} Creación de
  cuentas, cambio de rol y activación/desactivación de usuarios.
  \item \textbf{Mapa interactivo, geolocalización y rutas.} Mapa basado
  en OpenStreetMap (\texttt{flutter\_map}), posición actual del turista
  mediante el GPS del dispositivo (\texttt{geolocator}) y cálculo de la
  ruta por carretera hacia un destino con el servicio OSRM.
  \item \textbf{Clima en tiempo real.} Consulta de temperatura y estado
  del tiempo desde la API pública de Open-Meteo.
  \item \textbf{Reseñas y calificaciones.} Registro de reseñas con
  calificación por parte de los visitantes y visualización de las
  reseñas aprobadas de cada lugar, hotel, restaurante, etc.
  \item \textbf{Búsqueda y filtros.} Búsqueda general de contenido y
  filtro por categoría en cada módulo.
  \item \textbf{Versión web y móvil desde un mismo código.} La misma
  base de código Flutter genera la app Android y una versión web con
  diseño propio, desplegada en Netlify.
\end{enumerate}

\subsubsection*{Autenticación y sesión}

La aplicación no implementa un backend propio para la autenticación: se
apoya en Supabase Auth. Al iniciar sesión, el método
\texttt{signInWithPassword} devuelve una sesión compuesta por un
\emph{access token} (JWT, válido una hora con la configuración por
defecto de Supabase) y un \emph{refresh token}. El SDK
\texttt{supabase\_flutter} guarda la sesión en el dispositivo
(\texttt{SharedPreferences} en el móvil y \texttt{localStorage} en la
web) y renueva el JWT automáticamente antes de que venza, de modo que
el usuario permanece conectado hasta que cierra sesión.

La decisión de qué pantalla mostrar se centralizó en el widget
\texttt{AuthGate}, que escucha el flujo de cambios de estado de la
autenticación. Así, después de iniciar o cerrar sesión no es necesario
navegar manualmente: la interfaz reacciona sola al cambio
(Figura~\ref{fig:impl-authgate}).

\begin{figure}[H]
\begin{lstlisting}[style=codigo,language=Dart]
class AuthGate extends StatelessWidget {
  const AuthGate({super.key});

  @override
  Widget build(BuildContext context) {
    final client = Supabase.instance.client;

    return StreamBuilder<AuthState>(
      stream: client.auth.onAuthStateChange,
      builder: (context, snapshot) {
        if (client.auth.currentSession == null) {
          return const LoginScreen();
        }
        return const AdminHomeScreen();
      },
    );
  }
}
\end{lstlisting}
\caption{Selección de la pantalla según el estado de la sesión.
Fuente: \texttt{lib/screens/auth\_gate.dart}.}
\label{fig:impl-authgate}
\end{figure}

\subsubsection*{Autorización por rol en la base de datos}

Una decisión técnica central fue \textbf{no confiar la seguridad a la
interfaz}. Ocultar un botón en la app no impide que alguien llame
directamente a la API REST con la clave pública. Por eso los permisos
se definieron en PostgreSQL con \emph{Row Level Security} (RLS): cada
tabla tiene políticas que consultan las funciones
\texttt{es\_admin\_o\_editor()} y \texttt{es\_administrador()}, las
cuales identifican al usuario a partir de \texttt{auth.uid()} (el
\texttt{id} contenido en el JWT) y verifican además que la cuenta esté
activa (Figura~\ref{fig:impl-rls}). Los visitantes anónimos solo pueden
leer contenido activo; insertar o modificar requiere sesión con rol
válido.

\begin{figure}[H]
\begin{lstlisting}[style=codigo,language=SQL]
create or replace function public.es_admin_o_editor()
returns boolean
language sql stable
security definer
set search_path = public
as $$
  select exists (
    select 1 from public.usuarios u
    where u.id = auth.uid()
      and u.rol in ('administrador','editor') and u.activo
  );
$$;

create policy "lugares_insert_admin"
on public.lugares for insert
to authenticated
with check (public.es_admin_o_editor());
\end{lstlisting}
\caption{Función de rol y política RLS que restringe la inserción de
lugares turísticos. Fuente: \texttt{supabase/04\_SCHEMA\_TURISMO\_COMARAPA.sql}.}
\label{fig:impl-rls}
\end{figure}

Complementariamente, un \emph{trigger} sobre \texttt{auth.users}
(\texttt{on\_auth\_user\_created}) crea automáticamente el perfil en la
tabla \texttt{usuarios} con rol \texttt{editor} cada vez que se registra
una cuenta, y otro \emph{trigger} (\texttt{trg\_usuarios\_bloquear\_rol})
impide que un usuario que no es administrador cambie su propio rol desde
la aplicación.

\subsubsection*{Servicios externos: rutas y clima}

Para el cálculo de rutas se eligió el servidor público de OSRM, que no
requiere clave de API. El servicio \texttt{RoutingService} construye la
petición con las coordenadas de origen (GPS del turista) y destino, y
convierte la geometría GeoJSON de la respuesta en una lista de puntos
que se dibuja como polilínea sobre el mapa, junto con la distancia y el
tiempo estimado (Figura~\ref{fig:impl-osrm}). El clima se obtiene de
forma análoga desde Open-Meteo. En ambos casos se fijó un tiempo máximo
de espera (15 y 12 segundos) para que una red lenta no deje la pantalla
bloqueada.

\begin{figure}[H]
\begin{lstlisting}[style=codigo,language=Dart]
final uri = Uri.parse(
  'https://router.project-osrm.org/route/v1/driving/'
  '${origen.longitude},${origen.latitude};'
  '${destino.longitude},${destino.latitude}'
  '?overview=full&geometries=geojson',
);

final response = await http.get(uri).timeout(const Duration(seconds: 15));
if (response.statusCode != 200) {
  throw Exception('HTTP ${response.statusCode}');
}

final route = (jsonDecode(response.body)['routes'] as List).first;
final puntos = (route['geometry']['coordinates'] as List).map((c) {
  final par = c as List;
  return LatLng((par[1] as num).toDouble(), (par[0] as num).toDouble());
}).toList();
\end{lstlisting}
\caption{Cálculo de la ruta por carretera hacia un destino con OSRM
(fragmento simplificado). Fuente: \texttt{lib/services/routing\_service.dart}.}
\label{fig:impl-osrm}
\end{figure}

\subsubsection*{Dificultades encontradas y su resolución}

\begin{enumerate}
  \item \textbf{Crear un usuario cerraba la sesión del administrador.}
  Al llamar a \texttt{signUp} con el cliente principal de Supabase, la
  sesión del nuevo usuario reemplazaba a la del administrador que lo
  estaba creando. Se resolvió usando un cliente de Supabase
  \emph{temporal}, sin persistencia de sesión y con flujo
  \texttt{implicit}, que se descarta inmediatamente después del alta
  (Figura~\ref{fig:impl-tempclient}). Si el rol elegido es
  administrador, el ascenso se hace después con el cliente principal.

\begin{figure}[H]
\begin{lstlisting}[style=codigo,language=Dart]
final tempClient = SupabaseClient(
  config.supabaseUrl,
  config.supabasePublishableKey,
  authOptions: const AuthClientOptions(
    authFlowType: AuthFlowType.implicit,
  ),
);

try {
  final response = await tempClient.auth.signUp(
    email: correoCtrl.text.trim(),
    password: passwordCtrl.text,
    data: {'nombre': nombreCtrl.text.trim()},
  );
  newUserId = response.user?.id;
} finally {
  await tempClient.dispose();
}
\end{lstlisting}
\caption{Alta de un usuario con un cliente aislado para no reemplazar la
sesión del administrador (fragmento simplificado). Fuente:
\texttt{lib/screens/admin\_users\_screen.dart}.}
\label{fig:impl-tempclient}
\end{figure}

  \item \textbf{Error ``new row violates row-level security policy''
  al guardar contenido.} Las cuentas creadas antes de instalar el
  \emph{trigger} no tenían fila en \texttt{usuarios}, por lo que
  \texttt{es\_admin\_o\_editor()} devolvía falso y la base de datos
  rechazaba el guardado. Se preparó un script de diagnóstico que detecta
  las cuentas sin perfil y las completa
  (\texttt{supabase/05\_DIAGNOSTICO\_Y\_FIX\_USUARIOS.sql}).

  \item \textbf{Recursividad en las políticas de \texttt{usuarios}.} Una
  política de la tabla \texttt{usuarios} que consulta la misma tabla
  \texttt{usuarios} provoca una evaluación recursiva. Se evitó declarando
  las funciones de rol como \texttt{security definer} con
  \texttt{search\_path} fijo, de modo que se ejecutan sin volver a
  aplicar las políticas.

  \item \textbf{Las reseñas no se veían después de enviarlas.} El
  esquema inicial guardaba toda reseña como \texttt{pendiente} y la app
  solo muestra las \texttt{aprobada}, de modo que el visitante no veía
  su propia reseña. Se optó por la auto-aprobación: la política de
  inserción solo acepta filas marcadas como \texttt{aprobada}
  (\texttt{supabase/06\_RESENAS\_AUTOAPROBAR.sql}), sin abrir la puerta
  a otros estados.

  \item \textbf{Una sola base de código para móvil y web.} Las pantallas
  móviles no se veían bien en pantallas anchas, pero reescribir toda la
  navegación era costoso. Se creó una capa web propia
  (\texttt{lib/web}) que se elige con \texttt{kIsWeb} al iniciar la app,
  y un \texttt{NavigatorObserver} (\texttt{WebAuthRedirectObserver}) que,
  cuando una pantalla compartida abre el login o el panel móvil, los
  reemplaza por su versión web sin modificar esas pantallas.

  \item \textbf{Rutas de la versión web devolvían 404 en Netlify.} Al
  recargar una página distinta de la raíz, Netlify buscaba un archivo
  que no existe. Se agregó el archivo \texttt{web/\_redirects} con la
  regla \texttt{/* /index.html 200}, que entrega siempre la aplicación y
  deja que Flutter resuelva la ruta.

  \item \textbf{Falta de conexión o tablas vacías.} En la demostración
  no siempre había conexión ni datos cargados. Las pantallas de eventos,
  actividades y gastronomía cargan datos representativos de Comarapa
  cuando la consulta a Supabase falla o no devuelve registros, para que
  la interfaz nunca quede vacía.
\end{enumerate}
